\section{Comparison with Related Systems}
\label{sec:comparison}

Five open projects were read in full against this system, and mechanisms were
taken from each where they cost no model call. This section says what each
is, where the designs differ, and what was measured side by side.

\textbf{Codex and Open Math Harness.} Codex \cite{openaicodex2026} is a
coding agent whose turn loop assumes a frontier model with native tool calls;
Open Math Harness \cite{openmathharness2026} is a capability microkernel for
mathematics research. Neither was embedded. From Codex we took the patch
language with its lenient parser and fuzzy matcher, the append-only history,
grammar-constrained patches, truncation on insertion, budget reminders and
retry backoff; from Open Math Harness the Lean protocol, the scoring of
partial results and the candidate tree. What changed for small models on a
CPU is in Section~\ref{sec:harness}: compaction is written by the harness
rather than by the model, the grammar is rebuilt per step and leaves the
reasoning free, and the edit format is chosen by file length.

\textbf{vLLM Semantic Router} \cite{semanticrouter2026} is a decision layer
in front of many models: about twenty signal types, Boolean rules over them
that pick a decision, and a selection algorithm inside it. Its decision rule
is lexicographic and hand-weighted; ours is one rule derived from expected
loss with a price of latency in real units. It decides correctness by model
confidence, self-rating or peer agreement; ours by exact verifiers. Its
cascade serves the last response when none passes a threshold; ours returns
the best-scoring candidate or abstains. Its semantic cache stores every
successful response; our memory stores verified answers only. It reports its
own gate at 79\,ms to 9.8\,s per request on a CPU, because the signals come
from a model of 0.3B parameters or more; ours costs about 0.1\,s from
features the draft already produced. From it we took certified per-bin risk,
a declared policy for unknown features, the reask signal and the polarity
guard; and from its retrieval and memory code the model-free parts of a
conversation memory (the notes, the gate that decides when to read them, the
elbow threshold, duplicates that respect numbers) and the idea of a
topic-continuity signal, reduced to one rule for follow-up questions. Its
document search re-ranks the vector hits with BM25 and character trigrams
computed over those hits only; ours runs words and meaning as independent
arms over the whole base, which is what finds a passage the embedder missed.
Its check of an answer against its context is a token classifier of 150 to
300\,M parameters; ours is a rule over words and numbers
(Section~\ref{sec:kb:chat}). Its selection algorithms choose among many
models (Elo, nearest labelled queries, Thompson sampling, session-aware
switching); with two models there is nothing to select, and the one
question-only signal we could test did not add to the draft's own
(Table~\ref{tab:negative}). Its repository commits speed records and no
quality measurement for retrieval, memory, cache precision or selection. It
remains ahead on breadth: Envoy and Kubernetes integration, trained jailbreak
and personal-data classifiers, per-user memory isolation, a management
dashboard. Those need a classifier model per request or a service mesh, and
are outside what a single-binary assistant on four cores should carry.

\textbf{Memorilla} \cite{memorilla2026} trains a small module that reads a
document collection into sixteen vectors for a frozen decoder, on GPUs and
for decoders of 8B parameters. Our port has forward-pass parity with it and
serves its vectors through our own engine; what it cannot have on a CPU is
the pre-training (3.3 million rows in two stages). The comparison of the two
code bases explained our one failed module: a learning rate three times the
reference's at a quarter of its batch, no dropout, and checkpoint selection
by validation loss, which rewarded a module that had learnt a constant. The
reference has no test for that failure; the collapse statistics and the
activation guard of Section~\ref{sec:mem:tokens} are ours. Memorilla has no
run-time memory management (no write, forget or consolidation); that role is
taken here by the note memory of Section~\ref{sec:mem:recall}.

\textbf{DeepSeek Harness} \cite{deepseekharness2026} is a ReAct loop over an
event-sourced session log from which any prompt can be rebuilt byte for byte.
A run ends on the model's own claim; it has no verifier, no search over
candidates, no rollback and no step limit, and its sandbox does not unshare
the network. From it we took the anchored token meter, pruning before
compaction, the refusal of an edit from a truncated reply and the
repeat-attempt notice. It remains ahead on byte-exact prompt replay,
subagents at scale and browser tools.

\textbf{OntoSmith} \cite{ontosmith2026} builds a governed knowledge graph
under a human-approved, pinned ontology with deterministic gates and a
published precision record. Its served retrieval path has no BM25, fixed
windows without overlap and no rank fusion, where ours had all three. What it
does better, and what we took, is the role of the graph: the graph
\emph{supplies candidates} to the ranking instead of overriding it; an
ontology stored as data; and a quote for every relation.

\textbf{Measured side by side.} The benchmarks of Semantic Router that
compile offline were run interleaved with Rust counterparts, five rounds,
median (Table~\ref{tab:h2h}). Two results go against a simple reading. Their
hand-written AVX2 assembly is 1.3 to 1.5 times faster than our intrinsics for
long dot products; the difference is 50 to 60\,ns on a vector scan that costs
18 to 22\,ms a query. And the per-request cost of both signal layers is
negligible beside one model call, so what separates the projects is how an
answer comes to be accepted, not these numbers.

\begin{table}[t]
\centering\scriptsize
\caption{Per-request work, five interleaved rounds (load 24 to 29). The routing row is not the same computation (their 30 keyword rules against our domain classifier); the compression rows are; the last row compares host overhead only, against the other project's own budget.}
\label{tab:h2h}
\begin{tabular}{@{}p{3.3cm}rrr@{}}
\toprule
Workload & Other & Nafs-Aql & Ratio \\
\midrule
Routing signal, 2\,KB prompt & 197.7\,\textmu s & 2.2\,\textmu s & 90$\times$ \\
Extractive compression, 500 tok. & 745\,\textmu s & 310\,\textmu s & 2.4$\times$ \\
Extractive compression, 2K tok. & 5.07\,ms & 1.16\,ms & 4.4$\times$ \\
Extractive compression, 16K tok. & 50.3\,ms & 8.96\,ms & 5.6$\times$ \\
Dot product, $n=768$, AVX2 & 99\,ns & 145\,ns & 0.68$\times$ \\
Dot product, $n=1024$, AVX2 & 135\,ns & 199\,ns & 0.68$\times$ \\
Dot product, $n=1024$, portable & 1.20\,\textmu s & 1.21\,\textmu s & 0.99$\times$ \\
Agent host overhead, 100 turns & budget 900\,ms & 1.75\,ms & -- \\
\bottomrule
\end{tabular}
\end{table}
